\section{第11章の解答}\label{解答:第11章}

\paragraph{11.1}
$\sin\frac{4\pi}3=-\frac{\sqrt3}2$であり,\,$\arcsin$が返す角は$\left[-\frac\pi2,\frac\pi2\right]$にあるので
\[
\arcsin\left(\sin\frac{4\pi}3\right)=-\frac\pi3.
\]
$\cos\frac{5\pi}4=-\frac{\sqrt2}2$であり,\,$\arccos$が返す角は$[0,\pi]$にあるので
\[
\arccos\left(\cos\frac{5\pi}4\right)=\frac{3\pi}4.
\]
また$-\frac{2\pi}3$に$\pi$を加えると$\frac\pi3$になり,正接の値は変わらない.この角は$\left(-\frac\pi2,\frac\pi2\right)$の中にあるから
\[
\arctan\left(\tan\left(-\frac{2\pi}3\right)\right)=\frac\pi3.
\]
微分では合成関数の公式を使う.\,$-1<x<1$なら$0\le x^2<1$なので,逆正弦の微分の分母は正であり,
\[
\frac{\dd}{\dd x}\arcsin(x^2)
=\frac{2x}{\sqrt{1-(x^2)^2}}
=\frac{2x}{\sqrt{1-x^4}}
\]
となる.

\paragraph{11.2}
最初の分母を平方完成すると
\[
x^2+2x+5=(x+1)^2+4.
\]
$u=\frac{x+1}{2}$と置けば$\dd x=2\dd u$だから
\[
\int\frac1{x^2+2x+5}\dd x
=\frac12\int\frac1{1+u^2}\dd u
=\frac12\arctan\frac{x+1}{2}+C.
\]
分母は実数全体で正なので,この答えは実数全体で用いることができる.

2つ目の積分は,\,$\arcsin\frac x2$の微分が$\frac1{\sqrt{4-x^2}}$であることから
\[
\int_0^1\frac1{\sqrt{4-x^2}}\dd x
=\left[\arcsin\frac x2\right]_0^1
=\frac\pi6.
\]
三角置換$x=2\sin\theta$を使っても,\,$0\le\theta\le\frac\pi6$で分母と$\dd x$の係数がともに$2\cos\theta$となり,同じ値を得る.

\paragraph{11.3}
$\sin x\ne0$のところで商を微分すると
\[
(\cot x)'
=\left(\frac{\cos x}{\sin x}\right)'
=\frac{-\sin^2x-\cos^2x}{\sin^2x}
=-\frac1{\sin^2x}=-\csc^2x.
\]
さらに$\sin x\cos x\ne0$なら
\[
\tan x+\cot x
=\frac{\sin^2x+\cos^2x}{\sin x\cos x}
=\frac1{\sin x\cos x}.
\]
それぞれ分母の微分を使って積分すると
\[
\begin{aligned}
\int\frac1{\sin x\cos x}\dd x
&=\int\tan x\dd x+\int\cot x\dd x\\
&=-\log|\cos x|+\log|\sin x|+C\\
&=\log|\tan x|+C.
\end{aligned}
\]
もとの分母は$x=\frac{k\pi}{2}$で$0$になる.したがって,整数$k$について$\left(\frac{k\pi}{2},\frac{(k+1)\pi}{2}\right)$のような,これらの点を含まない1つの区間で用いる.任意定数は区間ごとに独立である.

\paragraph{11.4}
$2+\sin x\ge1$なので,元の被積分関数は$[0,2\pi]$全体で連続である.ただし$t=\tan\frac x2$は$x=\pi$で定義されない.そこで
\[
I_1=\int_0^\pi\frac1{2+\sin x}\dd x,\qquad
I_2=\int_\pi^{2\pi}\frac1{2+\sin x}\dd x
\]
に分ける.\,$x=\pi$から少し離れた区間で置換してから極限を取ると,本文と同じ変形によって
\[
I_1=\lim_{R\to\infty}\int_0^R\frac1{t^2+t+1}\dd t,\qquad
I_2=\lim_{S\to\infty}\int_{-S}^0\frac1{t^2+t+1}\dd t
\]
となる.\,$x\to\pi-$では$t\to+\infty$,\,$x\to\pi+$では$t\to-\infty$だからである.2つの極限を独立に求める.

原始関数の1つは
\[
F(t)=\frac2{\sqrt3}\arctan\frac{2t+1}{\sqrt3}.
\]
$\arctan t$の極限は$t\to+\infty$で$\frac\pi2$,\,$t\to-\infty$で$-\frac\pi2$だから
\[
\begin{aligned}
I_1&=\frac2{\sqrt3}\left(\frac\pi2-\frac\pi6\right)
=\frac{2\pi}{3\sqrt3},\\
I_2&=\frac2{\sqrt3}\left(\frac\pi6+\frac\pi2\right)
=\frac{4\pi}{3\sqrt3}.
\end{aligned}
\]
したがって求める値は$I_1+I_2=\frac{2\pi}{\sqrt3}$である.無限の端点が現れたのは置換によるものであり,元の積分に特異点があるわけではない.

\paragraph{11.5}
固定点$(-1,0)$から与えられた点へ引く直線の傾きは
\[
t=\cfrac{\cfrac5{13}}{1+\cfrac{12}{13}}
=\frac5{25}=\frac15.
\]
$m=5,n=1$とすれば
\[
m^2-n^2=24,\qquad2mn=10,\qquad m^2+n^2=26.
\]
3つの整数の最大公約数$2$で割って,\,$(12,5,13)$を得る.実際$12^2+5^2=13^2$であり,斜辺で割ると元の円上の点へ戻る.

$m,n$がともに奇数なら,\,$m^2,n^2$もともに奇数である.したがって差$m^2-n^2$と和$m^2+n^2$は偶数であり,\,$2mn$も偶数である.\,$t=\frac nm$を既約分数にしても,公式から得る3つの整数に共通因数がないとは限らない.

\paragraph{11.6}
定義を使って積を計算すると
\[
2\sinh u\cosh u
=2\frac{\e^u-\e^{-u}}2\frac{\e^u+\e^{-u}}2
=\frac{\e^{2u}-\e^{-2u}}2
=\sinh(2u).
\]
また$\cosh^2u=\frac12\{1+\cosh(2u)\}$なので
\[
\int_0^{\log2}\cosh^2u\dd u
=\left[\frac u2+\frac14\sinh(2u)\right]_0^{\log2}.
\]
$\e^{2\log2}=4$,\,$\e^{-2\log2}=\frac14$より
\[
\sinh(2\log2)=\cfrac{4-\cfrac14}{2}=\frac{15}{8}.
\]
したがって積分値は$\frac12\log2+\frac{15}{32}$である.

\paragraph{11.7}
$u>0$で$\cosh u$は$1$から無限大へ狭義単調増加するので,\,$x>1$と一対一に対応する.この範囲なら$\sinh u>0$であり,
\[
x=\cosh u,\qquad
\sqrt{x^2-1}=\sqrt{\sinh^2u}=\sinh u,\qquad
\dd x=\sinh u\dd u.
\]
したがって積分は$\int1\dd u=u+C$になる.元へ戻すには
\[
x+\sqrt{x^2-1}=\cosh u+\sinh u=\e^u
\]
を使えばよい.右辺は正なので
\[
\int\frac1{\sqrt{x^2-1}}\dd x
=\log\bigl(x+\sqrt{x^2-1}\bigr)+C\qquad(x>1).
\]
この問題では$u>0$の枝を選んだため,平方根を$\sinh u$と書けた.\,$u<0$の枝なら平方根は$-\sinh u$となり,同じ途中式をそのまま使うことはできない.

\paragraph{11.8}
まず$x=2\sinh u$と置く.\,$x=0$は$u=0$,\,$x=2$は$u=\operatorname{arsinh}1=\log(1+\sqrt2)$へ移る.後者を$U$と書けば
\[
\sinh U=1,\qquad\cosh U=\sqrt2,\qquad\sinh(2U)=2\sqrt2.
\]
$\sqrt{4+x^2}=2\cosh u$,\,$\dd x=2\cosh u\dd u$だから
\[
\begin{aligned}
\int_0^2\sqrt{4+x^2}\dd x
&=4\int_0^U\cosh^2u\dd u\\
&=\left[2u+\sinh(2u)\right]_0^U\\
&=2\log(1+\sqrt2)+2\sqrt2.
\end{aligned}
\]

次に$t=x+\sqrt{4+x^2}$と置く.和と差の積が$4$なので
\[
\sqrt{4+x^2}-x=\frac4t.
\]
和と差を取って
\[
x=\frac12\left(t-\frac4t\right),\qquad
\sqrt{4+x^2}=\frac12\left(t+\frac4t\right),\qquad
\dd x=\frac12\left(1+\frac4{t^2}\right)\dd t
\]
を得る.端点は$t=2$と$t=2+2\sqrt2$である.したがって
\[
\begin{aligned}
\int_0^2\sqrt{4+x^2}\dd x
&=\frac14\int_2^{2+2\sqrt2}
 \left(t+\frac8t+\frac{16}{t^3}\right)\dd t\\
&=\left[\frac{t^2}{8}+2\log t-\frac2{t^2}\right]_2^{2+2\sqrt2}.
\end{aligned}
\]
べきの項を直接代入してもよいが,
\[
\frac12x\sqrt{4+x^2}=\frac{t^2}{8}-\frac2{t^2}
\]
と見れば,上端で$2\sqrt2$,下端で$0$になる.対数の差は
\[
2\log(2+2\sqrt2)-2\log2=2\log(1+\sqrt2)
\]
なので,1つ目の方法と一致する.さらに$x=2\sinh u$を代入すれば
\[
t=2\sinh u+2\cosh u=2\e^u
\]
となる.2つの置換は,変数をこの式で換えたものになっている.
